\section{What did not hold}
\label{sec:nt-failures}

The locked plan was titled \emph{``Branded by People, Polished by Machines: The 2024 Naming Turn in ICLR Titles\ldots''}. It predicted a dated 2024 break, communal change within incumbents, an LLM register that turns over with each model generation, and a review stage that depletes LLM markers. Several of these predictions were set from pilot estimates on the same rows, and they failed anyway. Table~\ref{tab:hypotheses} gives every verdict with its pilot value, and Appendix~\ref{app:nt-failures} lists the failures. In brief: no dated 2024 break (H1a); content is not a minority (H1b); unexplained drift did not demonstrably accelerate (H1c); the tool share is not a bound (H1e); communal change is not identified (H2a); newcomer-free teams changed less than predicted (H2c); no register turnover or vintage alignment (H3b, H3c); the lexical-over-structural ratio is short of $1.5\times$ (H3d); and review does not deplete LLM markers or understate LLM use as predicted (H4b, H4c).

\paragraph{Why this matters.}
Three failures are failures of \emph{identification}, and each generalises. \textbf{Shared outcomes}: the communal-change statistic meets its threshold even with no communal change, because co-authors share one title. \textbf{Collider bias}: the same-paper comparator for the lexical signature conditions on a common effect of both variables. \textbf{Ten cohorts}: paper-level standard errors made a year-level treatment look decisive (Holm $p$ 0.0006 $\rightarrow$ 0.071 for H1a). The adversarial reviews moved numbers in both directions: a calibration bug had \emph{inflated} the content share, and a small first audit had made \NAME{} neutrality look fragile until the complete audit made it robust. A gate without multiplicity control would have let through two acceptance associations that disappear under Holm.
